% Pattern: number/e-companion. \ECSwitch of an INFORMS class starts the
% e-companion, which numbers the theorems again from one after EC. Seen in the
% gold pilot (set X).
% want theorem-1 statement/theorem 1 status=none
% want lemma-1 statement/lemma 1 status=none
% want theorem-2 statement/theorem 2 status=none
% want lemma-2 statement/lemma EC.1 status=none
% want lemma-3 statement/lemma EC.2 status=none
% want theorem-3 statement/theorem EC.1 status=none
\documentclass{informs3}
\newtheorem{theorem}{Theorem}
\newtheorem{lemma}{Lemma}
\begin{document}
\section{Introduction}
\begin{theorem} Every compact space is bounded. \end{theorem}
\begin{lemma} Every closed subset of a compact space is compact. \end{lemma}
\section{Products}
\begin{theorem} Every finite product of compact spaces is compact. \end{theorem}
\ECSwitch
\section{Proofs}
\begin{lemma} Every compact metric space is complete. \end{lemma}
\begin{lemma} Every compact metric space is separable. \end{lemma}
\begin{theorem} Every product of compact spaces is compact. \end{theorem}
\end{document}
